\documentclass[10pt,a4paper]{amsart}
\usepackage[margin=19mm]{geometry}
\usepackage{amsmath,amssymb,mathtools}
\usepackage[scr=rsfs,cal=euler]{mathalfa}
\usepackage{xfrac}
\usepackage[unicode,hidelinks,bookmarks=false]{hyperref}
\newcommand{\ie}{i.e.\ }
\newcommand{\cf}{cf.\ }
\newcommand{\resp}{respectively\ }
\newcommand{\Subsection}{\subsection}
\providecommand{\nobreakdash}{\nobreak\hbox{-}}
\setlength{\emergencystretch}{2em}
\multlinegap0pt
\usepackage{bm}
\usepackage{bbm}
\usepackage{dsfont}
\usepackage{xspace}
\usepackage{cancel}
\usepackage{mathtools}
\usepackage{amsthm}
\makeatletter
\@ifundefined{theorem}{\newtheorem{theorem}{Theorem}}{}
\@ifundefined{lemma}{\newtheorem{lemma}[theorem]{Lemma}}{}
\@ifundefined{proposition}{\newtheorem{proposition}{Proposition}}{}
\makeatother
\title[Paperfolding Structures as Templates for Horseshoes in Multi]{Paperfolding Structures as Templates for Horseshoes in Multidimensional H\'enon Maps}
\author{Jizhou Li, Keisuke Fujioka, Akira Shudo}
\date{Source: arXiv:2509.08251v1 (CC BY 4.0)}
\begin{document}
\maketitle

\noindent\emph{Source and licence.} Paperfolding Structures as Templates for Horseshoes in Multidimensional H\'enon Maps. Version arXiv:2509.08251v1 (CC BY 4.0), distributed under the Creative Commons Attribution 4.0 International licence, as recorded in the arXiv licence metadata for this exact version and shown on the abstract page https://arxiv.org/abs/2509.08251v1. Full rights notice: licenses/arxiv-2509.08251-CC-BY-4.0.txt.\par\smallskip

\noindent\textit{Editorial scope.} Counted unit: the Proposition whose source label is \texttt{Semi-conjugacy to the full shift on N symbols} in the article (source0.tex lines 187--203, source label \texttt{Semi-conjugacy to the full shift on N symbols}), delivered together with the article's own complete original proof of it (source0.tex lines 205--271, one \texttt{proof} environment of 67 lines). No mathematics is written, repaired or re-derived: statement and proof are the article's own text line for line. Required context retained uncounted: eq:R definition (lines 122--124). Every definition, display and label of those lines is retained unchanged. Omitted from this package and not redistributed: 1--121, 125--186, 204--204, 272--1801. Nothing inside a retained excerpt is omitted or altered: every retained body is delivered exactly as the article's lines are written, so no omission receipt of any kind accompanies this package. Citations: the retained excerpts carry no citation command at all, so no bibliography entry of the article is transcribed and no reference is redistributed. Reference adaptation: none was needed and none was made; every reference inside the delivered unit resolves inside it, so no unresolved reference or citation is produced. Printed slips kept literal: no printed word, symbol, hypothesis or number of the counted statement or of its proof was altered. Layout and class: the article's own class is \texttt{iopart} with packages including graphicx, subcaption, xcolor, cite, amsthm, amssymb, inputenc; the wrapper is the lane's pinned \texttt{amsart} editorial wrapper at 10pt on A4, which restates the theorem-like environments the retained text needs and adds the article's own macro definitions that the retained text needs (none), with extra wrapper packages (bm, bbm, dsfont, xspace, cancel, mathtools). The delivered numbering is the unit's own. Assets: no retained span contains a figure environment, a graphics-inclusion command, a PDF-page inclusion, a table or a TikZ picture; no image file is copied into this package, the package declares no asset dependency and no image file is redistributed with it. Licence: version arXiv:2509.08251v1 of this article is distributed under the Creative Commons Attribution 4.0 International licence, as recorded in the arXiv licence metadata for this exact version and shown on the abstract page https://arxiv.org/abs/2509.08251v1. A full rights notice is delivered at licenses/arxiv-2509.08251-CC-BY-4.0.txt. Characters: every retained excerpt is pure ASCII, so the delivered engine, XeLaTeX with its default Latin Modern text font, reports no missing character.
\begin{equation}\label{eq:R definition}
R = I^u \times I^s \subset \mathbb{R}^n,
\end{equation}

\begin{proposition}[Semi-conjugacy to the full shift on \( N \) symbols]\label{Semi-conjugacy to the full shift on N symbols}
Let \( f : \mathbb{R}^n \to \mathbb{R}^n \) be a homeomorphism, and let \( R = I^u \times I^s \subset \mathbb{R}^{n_u} \times \mathbb{R}^{n_s} \) be the compact hypercylinder defined in Eq.~(\ref{eq:R definition}). Suppose that for every \( s \in I^s \), the horizontal disk
\[
d(s) := I^u \times \{s\} \subset R
\]
satisfies
\[
f(d(s)) \cap R = h_1(s) \sqcup h_2(s) \sqcup \dots \sqcup h_N(s),
\]
where each \( h_i(s) \subset R \) is a horizontal slice of \( R \).

Then:
\begin{enumerate}
    \item There exists a nonempty compact invariant set \( \Lambda \subset \bigcap_{n \in \mathbb{Z}} f^n(R) \).
    \item The restricted map \( f|_\Lambda \) is semi-conjugate to the full shift on \( N \) symbols.
\end{enumerate}
\end{proposition}

\begin{proof}
For each \( i = 1, \dots, N \), define the horizontal slab
\[
H_i := \bigcup_{s \in I^s} h_i(s) \subset R.
\]
Each \( H_i \) is a horizontal slab of \( R \), as it consists of horizontal slices varying continuously in \( s \), and the collection \( \{H_i\}_{i=1}^N \) is pairwise disjoint.

We now show that the preimage \( V_i := f^{-1}(H_i) \subset R \) is a vertical slab. For each \( s \in I^s \), since \( h_i(s) \subset f(d(s)) \), we have \( f^{-1}(h_i(s)) \subset d(s) \). Each preimage is compact because \( f \) is a homeomorphism and \( h_i(s) \) is compact. As \( s \) varies, the slices \( h_i(s) \) vary continuously in the Hausdorff metric, and so do their preimages. Define
\[
V_i := \bigcup_{s \in I^s} f^{-1}(h_i(s)).
\]
Then \( V_i \) is a vertical slab of \( R \).

Let \( \Sigma_N := \{1, \dots, N\}^\mathbb{Z} \) be the full shift space. For each finite forward sequence \( (i_0, i_1, \dots, i_k) \), define
\[
V_{i_0, i_1, \dots, i_k} := V_{i_0} \cap f^{-1}(V_{i_1}) \cap \dots \cap f^{-k}(V_{i_k}).
\]
We claim that each \( V_{i_0, \dots, i_k} \) is a vertical slab and nonempty. This follows inductively from the following lemma.

\begin{lemma}
Let \( V \subset V_i \) be a vertical slab and \( a \in \{1, \dots, N\} \). Then \( f^{-1}(V) \cap V_a \) is a vertical slab of \( R \).
\end{lemma}

\begin{proof}
Define
\[
K_a(s) := f^{-1}(V \cap h_a(s)).
\]
Since \( h_a(s) \subset f(d(s)) \), the preimage \( f^{-1}(V \cap h_a(s)) \subset d(s) \). Each \( K_a(s) \) is closed, and the map \( s \mapsto K_a(s) \) varies continuously in the Hausdorff topology. Define
\[
\tilde{V}_a := \bigcup_{s \in I^s} K_a(s) = f^{-1}(V) \cap V_a.
\]
Then \( \tilde{V}_a \subset R \) is a vertical slab.
\end{proof}

Using this lemma inductively, all forward intersections \( V_{i_0, \dots, i_k} \) are nonempty vertical slabs.

Now define, for any backward sequence \( (i_{-1}, \dots, i_{-k}) \),
\[
H_{i_{-1}, \dots, i_{-k}} := H_{i_{-1}} \cap f(H_{i_{-2}}) \cap \dots \cap f^{k-1}(H_{i_{-k}}).
\]
Similar derivations show that all such sets are nonempty horizontal slabs of \( R \).

Fix any bi-infinite sequence \( \mathbf{i} = (\dots, i_{-1}, i_0, i_1, \dots) \in \Sigma_N \). For each \( k \geq 0 \), define
\[
\Lambda_k := V_{i_0, \dots, i_k} \cap H_{i_{-1}, \dots, i_{-k}}.
\]
Each \( \Lambda_k \subset R \) is compact and nonempty, and the sequence \( \{\Lambda_k\} \) is nested. Define
\[
\Lambda_{\mathbf{i}} := \bigcap_{k \geq 0} \Lambda_k.
\]
Then \( \Lambda_{\mathbf{i}} \subset R \) is a nonempty compact set. In general, \( \Lambda_{\mathbf{i}} \) may contain more than one point, since no assumptions are made about the widths of the slabs.

Define the invariant set
\[
\Lambda := \bigcup_{\mathbf{i} \in \Sigma_N} \Lambda_{\mathbf{i}} \subset \bigcap_{n \in \mathbb{Z}} f^n(R).
\]
The map
\[
\pi : \Lambda \to \Sigma_N, \quad \pi(x) = \mathbf{i}\quad \mathrm{ if }\quad x \in \Lambda_{\mathbf{i}}
\]
is well-defined and continuous. Moreover,
\[
\pi(f(x)) = \sigma(\pi(x)),
\]
where \( \sigma \) is the left shift on \( \Sigma_N \). Thus, \( f|_\Lambda \) is semi-conjugate to the full shift on \( N \) symbols.
\end{proof}

\end{document}
